\documentclass[11pt]{article}
\topmargin=-1.0cm
\usepackage{graphicx}
\usepackage{epstopdf}
\usepackage{amssymb}
\usepackage{amsmath}
\usepackage{amsthm}
\usepackage{color}
\usepackage[margin=1in]{geometry}
\usepackage[hidelinks]{hyperref}
\usepackage{booktabs}

% ******************************* Author Macros *******************************

\newcommand{\pl}[2]{\frac{\partial#1}{\partial#2}}
\newcommand{\bu}{{\bf u} }
\newcommand{\p}{\partial}
\newcommand{\og}{\omega}
\newcommand{\Og}{\Omega}
\newcommand{\fl}[2]{\frac{#1}{#2}}
\newcommand{\dt}{\delta}
\newcommand{\nn}{\nonumber}
\newcommand{\ap}{\alpha}
\newcommand{\bt}{\beta}
\newcommand{\tht}{\theta}
\newcommand{\kp}{\kappa}
\newcommand{\Dt}{\Delta}
\renewcommand{\theequation}{\arabic{section}.\arabic{equation}}
\newcommand{\be}{\begin{equation}}
\newcommand{\ee}{\end{equation}}
\newcommand{\ba}{\begin{array}}
\newcommand{\ea}{\end{array}}
\newcommand{\bea}{\begin{eqnarray}}
\newcommand{\eea}{\end{eqnarray}}
\newcommand{\beas}{\begin{eqnarray*}}
\newcommand{\eeas}{\end{eqnarray*}}
\newtheorem{theorem}{Theorem}[section]
\newtheorem{lemma}{Lemma}[section]
\newtheorem{proposition}{Proposition}[section]
\newtheorem{remark}{Remark}[section]
\newtheorem{assumption}{Assumption}[section]
\newcommand{\bx}{{\bf x} }
\newcommand{\gm}{\gamma}
\newcommand{\vep}{\varepsilon}
\newcommand{\T}{\mathbb T}
\newcommand{\dT}{d_{\mathbb T}}

% *****************************************************************************

\title{A WKB-based fixed-grid method for capturing trait concentration in a dispersal evolution model}

\author{
Weijie Huang\thanks{School of Mathematics and Statistics, Beijing Jiaotong University, Beijing 100044, People's Republic of China {\tt (wjhuang@bjtu.edu.cn)}} \ and
Xinran Ruan\thanks{School of Mathematical Sciences, Capital Normal University, Beijing, China, 100048, People's Republic of China {\tt (xinran.ruan@cnu.edu.cn)}.}
}

\begin{document}
\date{}
\maketitle

\begin{abstract}
The evolution of dispersal traits is a fundamental topic in evolutionary ecology, 
where natural selection may drive the trait distribution toward concentration in the rare mutation regime. 
This singular behavior poses a serious numerical difficulty, 
since direct discretizations of the population density require very fine trait grids to identify the fittest trait and to resolve the concentrated profile accurately. 
In this paper, we develop a WKB-based numerical framework for a dispersal evolution model. 
By separating the exponentially concentrated trait dependence from a smoother
amplitude variable and combining this WKB representation with specially
designed reconstruction techniques, the method recovers the selected trait and
the associated concentration structure accurately and efficiently on fixed trait grids. 
We establish a semi-discrete stationary fixed-grid asymptotic-preserving structure for the rare-mutation limit of the steady-state problem.
Numerical experiments compare the proposed method with
direct density discretizations and confirm its advantage in the small mutation
regime.
\end{abstract}

{\bf Key words. } dispersal evolution, rare mutation limit, trait concentration, WKB reformulation, Hamilton--Jacobi limit, asymptotic-preserving structure.

{\bf MSC 2020. } 65M06, 65M12, 35B25, 35F21, 92D15.


%===================================================================================================
%	Introduction
%===================================================================================================
\section{Introduction}

Adaptive dynamics studies the evolution of phenotypic traits under mutation, selection, and ecological feedback. In structured population models, natural selection is often manifested by the concentration of the population density on one or several fittest traits. In the rare-mutation regime, such concentration is usually described by a measure-valued limit in the trait variable, and the support of the limiting measure represents the selected phenotype. This viewpoint has been developed in selection-mutation models and in the Hamilton--Jacobi approach to adaptive dynamics \cite{BarlesPerthame, DiekmannJabinMischlerPerthame}. Among these problems, the evolution of dispersal is a particularly important example. In spatially heterogeneous environments, different dispersal rates may lead to different fitness advantages, and the selected strategy depends on the interaction between movement, resource distribution, and population competition. This question has been studied from several viewpoints in mathematical ecology, for example in reaction-diffusion models for slow dispersal, ideal free distributions, and spatial sorting \cite{DockeryHutsonMischaikowPernarowski1998, Hastings1983,Ronce2007}.

In this work, we consider a representative dispersal evolution model of the form 
\begin{equation}\label{eq:n_intro} 
	\varepsilon \partial_t n_\varepsilon -D(\theta)\Delta_{\mathbf{x}} n_\varepsilon -\varepsilon^2\Delta_\theta n_\varepsilon = n_\varepsilon\bigl(K(\mathbf{x})-\rho_\varepsilon(\mathbf{x})\bigr), \qquad \mathbf{x}\in\Omega . \end{equation} 
Here \(t\) denotes time, \(\mathbf{x}\) is the spatial variable, and \(\theta\) is the phenotypic trait. The unknown
\(n_\varepsilon(t,\mathbf{x},\theta)\) is the structured population density, and 
\[ 
	\rho_\varepsilon(t,\mathbf{x}) = \int_{\mathbb T} n_\varepsilon(t,\mathbf{x},\theta)\,d\theta, \quad 
	P_\varepsilon(t,\theta) = \int_\Omega n_\varepsilon(t,\mathbf{x},\theta)\,d\mathbf{x},  
\]
denote the population density in space and the trait marginal, respectively. 
The concentration of \(P_\varepsilon\) in the rare-mutation regime is used to identify the selected trait. 
The function \(K(\mathbf{x})\) represents the local carrying capacity, while the trait-dependent diffusivity \(D(\theta)\) models the dispersal rate associated with phenotype \(\theta\). Throughout the paper, we assume that the dispersal rate is positive and periodic in the trait variable and that the carrying capacity is nonnegative and bounded: 
 \begin{equation}\label{eq:basic_assumptions} 
 	0<D_{\min}\le D(\theta)\le D_{\max}<\infty, \qquad 0\le K(\mathbf{x})\le K_{\max}. 
\end{equation} 
We impose homogeneous Neumann boundary conditions on \(\partial\Omega\) and periodic boundary conditions in the trait direction. The time-dependent model \eqref{eq:n_intro} is used here as an evolutionary relaxation toward the selection state, while the asymptotic structure of interest is motivated by the rare-mutation limit of the associated steady-state problem.


The steady-state version of this dispersal evolution model was studied by
B. Perthame and P. E. Souganidis in \cite{PerthameSouganidis}.  Their work provides a
rare-mutation description of the selected dispersal trait for a population
structured by space and by a continuous dispersal trait.  In the limit of
vanishing mutations, the population concentrates in the trait variable, while
the spatial profile remains nontrivial.  This result gives the asymptotic
background for the present paper and clarifies the limiting selection structure
that the numerical method aims to capture.




This dispersal model fits into a broader class of structured selection models with concentration in the small-mutation limit. For purely trait-structured or nonlocal selection models, Dirac concentration and constrained Hamilton--Jacobi equations have been studied in \cite{DesvillettesJabinMischlerRaoul} and \cite{LorzMirrahimiPerthame}. When space is included as an additional structuring variable, the analysis becomes more delicate. 
Related Hamilton--Jacobi limits for populations structured by space and trait were obtained in \cite{BouinMirrahimi}, and an asymptotic analysis of a selection model with space was given in \cite{MirrahimiPerthame}. 
The Perthame--Souganidis dispersal model has also motivated later studies on dispersal evolution, conditional dispersal, and the stability of Dirac concentrations \cite{HaoLamLou2019, Lam2017,LamLouPerthame2023}.



The same concentration mechanism creates a severe numerical difficulty.  Direct
discretizations of the original density must resolve a trait profile whose
width shrinks with the mutation scale.  Hence the mesh required in the trait
direction may become prohibitively fine when \(\varepsilon\) is small.  This
suggests that the numerical method should exploit the asymptotic separation
between the singular trait dependence and the smoother amplitude, rather than
resolve the concentrated density directly.

This is closely related to the idea of asymptotic-preserving schemes, whose
goal is to capture the limiting regime on discretizations that do not resolve
the small scale \(\varepsilon\) \cite{JinAP}.  In rare-mutation problems, the
natural asymptotic separation is provided by a WKB or Hopf--Cole
transformation.  The density is represented by an exponential phase, which
describes the concentration in the trait variable, multiplied by a smoother
amplitude.  Such a WKB viewpoint is standard in the Hamilton--Jacobi approach
to adaptive dynamics \cite{BarlesPerthame} and is also central in the
dispersal model of \cite{PerthameSouganidis}.

WKB-type transformations have been used as a numerical tool in several multiscale problems with singular or highly oscillatory structures. In adaptive dynamics, a WKB representation was used in \cite{AlmeidaPerthameRuan} to construct an asymptotic-preserving scheme that captures Dirac concentrations on coarse, \(\varepsilon\)-independent phenotype meshes. Related Hopf--Cole based AP ideas have also been used for kinetic reaction-transport equations in singular front propagation regimes \cite{Hivert2018}. In kinetic BGK models, a Hopf--Cole transformation has been used to approximate the density through an effective phase equation and higher order correction terms \cite{LuoPayne2017}. WKB-based numerical ideas are also well established for highly oscillatory semiclassical problems, where a WKB-type preprocessing separates the dominant oscillations and allows the remaining smoother problem to be discretized on much coarser grids \cite{ArnoldAbdallahNegulescu2011,KoernerArnoldDoepfner2022}.



%The same concentration mechanism creates a severe numerical difficulty.  Direct
%discretizations of the original density must resolve a trait profile whose width
%shrinks with the mutation scale.  Hence the mesh required in the trait direction
%may become prohibitively fine when \(\varepsilon\) is small.  WKB-based
%asymptotic-preserving methods offer a natural way to avoid this resolution
%barrier.  In particular, for age- and trait-structured adaptive dynamics
%models, a WKB representation has been used to construct AP schemes that capture
%Dirac concentrations on coarse, \(\varepsilon\)-independent phenotype meshes
%\cite{AlmeidaPerthameRuan}.  This is consistent with the general philosophy of
%asymptotic-preserving schemes for multiscale problems \cite{JinAP}.  Numerical
%Hamilton--Jacobi discretizations, including monotone and high-order
%ENO/WENO-type schemes, provide the standard tools for approximating the phase
%equation \cite{CrandallLions,OsherShu,JiangPeng,ShuHJ}.

%The dispersal evolution problem considered here has an additional coupling
%structure.  The effective fitness is determined by a spatial eigenvalue problem
%depending on the reconstructed marginal density, while the selected trait is
%encoded in the maximum of the WKB phase.  Thus it is not enough to evolve the
%phase and amplitude variables.  One also has to reconstruct the density-level
%quantities in a stable way.  This is particularly important for the trait
%marginal, whose peak location, height, and width are all sensitive to small
%errors in the phase when \(\varepsilon\) is small.


%Model \eqref{eq:n_intro} and related equations have been extensively studied from the analytical viewpoint, including well-posedness, long-time behavior, and the connection with adaptive dynamics \cite{MirrahimiPerthame,PerthameSouganidis,AlmeidaPerthameRuan}. A characteristic feature is that the competition between spatial heterogeneity and trait-dependent dispersal may select an optimal phenotype in the long-time regime. In many situations, the steady state exhibits concentration in the trait variable. The small parameter $\varepsilon>0$ plays the role of a rare mutation scale, and in the limit $\varepsilon\to0$ the steady state typically converges, in a weak sense, to a Dirac mass supported at the fittest trait. This singular limit is well understood at the PDE level through WKB or Hopf--Cole transforms and constrained Hamilton--Jacobi equations coupled with principal eigenvalue problems \cite{PerthameSouganidis,BarlesPerthame,MirrahimiPerthame}.
%
%From the numerical point of view, however, this asymptotic regime remains challenging. As $\varepsilon$ decreases, the trait profile of $n_\varepsilon$ becomes increasingly narrow, so standard discretizations of the original variable require a rapidly growing number of grid points in the phenotypic direction. Moreover, approximating the full density with reasonable accuracy does not automatically guarantee an accurate identification of the dominant trait. This is particularly restrictive in the small-$\varepsilon$ regime, where the quantity of real interest is often the location of the concentration and its effective profile rather than a pointwise resolution of the entire density on an extremely fine mesh.
%
%Motivated by the asymptotic structure of the rare-mutation limit, we adopt the WKB representation
%\[
%n_\varepsilon(\mathbf{x},\theta,t)
%=
%W_\varepsilon(\mathbf{x},\theta,t)\,
%\exp\!\left(\frac{u_\varepsilon(\theta,t)}{\varepsilon}\right),
%\]
%where the phase variable $u_\varepsilon$ describes the rapidly varying exponential profile in the trait direction, while $W_\varepsilon$ is a comparatively smoother amplitude. In the long-time and rare-mutation regime, the effective behavior of the phase variable is related to a constrained Hamilton--Jacobi equation involving a principal eigenvalue problem in the spatial variable \cite{PerthameSouganidis,AlmeidaPerthameRuan}. This observation suggests that the WKB reformulation is not only asymptotically natural but also computationally advantageous. In particular, it provides a mechanism for extracting the dominant trait from the phase variable without directly resolving the narrow peak of the original density on a prohibitively fine trait grid.
%



The present paper develops a WKB-based fixed-grid numerical framework for the
dispersal evolution problem \eqref{eq:n_intro}.  Starting from the
phase-amplitude representation, we derive a semi-discrete scheme for the phase
and amplitude variables and prove the positivity of the amplitude.  We then
establish a stationary fixed-grid AP structure associated with the
rare-mutation limit of the steady-state problem under a one-sided remainder
stability condition.  
At the fully discrete level, the time-marching scheme is
coupled with a scalar mass-balance projection and a reconstruction layer.  The
reconstruction layer recovers the spatial marginal, the trait marginal, and the
local concentration structure from the WKB variables.  
It is essential for
accurate small-\(\varepsilon\) computations and also supports the dual
trait-grid implementation, in which the one-dimensional phase is resolved on a
finer trait grid while the space-dependent amplitude is evolved on a coarser
grid.


%The present paper is organized around this computational viewpoint. Our main goal is to develop and study a WKB-based numerical strategy that is able to capture trait concentration more efficiently than direct discretizations of the original model. To this end, we first derive a semi-discrete formulation for the phase and amplitude variables. We then specify a concrete numerical method based on a monotone discretization of the Hamilton--Jacobi part and a positivity-preserving discretization of the amplitude equation. At the analytical level, we establish positivity of the amplitude and derive a stationary, fixed-grid AP structure under transparent reconstruction and one-sided remainder stability assumptions. Since the WKB split does not automatically provide an exact discrete chain rule, we also record a density-compatible variant that restores a conservative density-level balance and makes the weighted remainder estimate verifiable. {\color{red} At the implementation level, a second theme is equally important: after the WKB variables have been advanced, the quantities used in the nonlocal reaction and in the diagnostics must be reconstructed carefully.  The current code therefore reconstructs \(\rho\) by log-stabilized quadrature or by a hybrid Laplace rule, reconstructs the trait marginal through \(\log P\) rather than by interpolating \(P\) itself, uses WENO interpolation of \(\log P\) to locate the peak, applies a scalar mass-balance correction, and fits a local quadratic profile to estimate the peak width.  These reconstruction steps are numerical ingredients of the method, not merely plotting devices, because they determine the robustness of small-\(\varepsilon\) computations.} On the numerical side, our experiments are designed to compare the original discretization and the WKB formulation directly, to track the remainder diagnostic, and to test the accurate capture and reconstruction of the fittest trait when $\varepsilon$ is small.
%





The paper is organized as follows.  Section~2 recalls the WKB ansatz and the
steady-state rare-mutation structure that motivates the numerical formulation.
Section~3 presents the semi-discrete WKB scheme, proves positivity of the
amplitude, and establishes the stationary fixed-grid AP structure.  Section~4
is devoted to the fully discrete method, including the scalar mass-balance
projection, the reconstruction of density-level quantities, and the dual
trait-grid implementation.  Section~5 contains numerical experiments comparing
the WKB method with direct density discretizations and testing the
reconstruction of the selected trait in the small-mutation regime.  The
appendices collect the auxiliary numerical ingredients and implementation
details used in the main text.



%The paper is organized as follows. In Section~2 we recall the WKB ansatz and the formal steady-state rare mutation limit that motivate the numerical formulation. In Section~3 we present the semi-discrete scheme, prove positivity, and derive the conditional stationary AP structure, including the density-compatible variant that controls the weighted remainder. {\color{red} We then add a separate subsection on the fully discrete reconstruction pipeline, where the phase and amplitude updates are stated together with the \(\rho\)- and \(P(\theta)\)-reconstructions, mass correction, residuals, time-step restrictions, and dual trait-grid option used in the code.} Section~4 is devoted to numerical experiments, where we compare the WKB formulation with direct discretizations, examine its performance in the concentration regime, and monitor the remainder diagnostic. The appendices collect typical choices of monotone numerical Hamiltonians and monotone numerical fluxes, density and trait-marginal reconstruction details, and sufficient conditions for the one-sided remainder bound.


%===================================================================================================
%   Preliminary results
%===================================================================================================
\section{Preliminary}

\subsection{WKB ansatz and reformulated system}

To capture the concentration phenomenon in the trait variable, we introduce the classical WKB ansatz
\begin{equation}\label{eq:WKB}
   n_\varepsilon(x,\theta,t)
   = W_\varepsilon(x,\theta,t)\,
   \exp\!\left(\frac{u_\varepsilon(\theta,t)}{\varepsilon}\right).
\end{equation}
This decomposition separates the exponentially varying part in the trait direction from a comparatively smoother amplitude. It is therefore natural for both the continuous asymptotic analysis and the numerical design in the small-mutation regime.

Substituting \eqref{eq:WKB} into \eqref{eq:n_intro} and assigning the leading trait-direction exponential contributions to the phase function, we obtain the reformulated system
\begin{equation}\label{eq:WKB_reformulation0}
\left\{
\begin{aligned}
& \partial_t u_\varepsilon
- |\partial_\theta u_\varepsilon|^2
- \varepsilon \, \partial_{\theta\theta} u_\varepsilon
= - \mathcal H(\theta,t), \\
& \varepsilon \partial_t W_\varepsilon
- D(\theta)\Delta_x W_\varepsilon
- \varepsilon^2 \partial_{\theta\theta} W_\varepsilon
- 2\varepsilon \, \partial_\theta W_\varepsilon \, \partial_\theta u_\varepsilon
= W_\varepsilon\bigl(K(x)-\rho_\varepsilon(x,t)+\mathcal H(\theta,t)\bigr),
\end{aligned}
\right.
\end{equation}
where \(\mathcal H(\theta,t)\) denotes an effective Hamiltonian. 
Using the identity
\[
\partial_\theta(W_\varepsilon\partial_\theta u_\varepsilon)
=
\partial_\theta W_\varepsilon\,\partial_\theta u_\varepsilon
+
W_\varepsilon\partial_{\theta\theta}u_\varepsilon,
\]
we rewrite the transport coupling in conservative form. If
\begin{equation}\label{def:F}
F_\varepsilon=-W_\varepsilon\partial_\theta u_\varepsilon,
\end{equation}
then
\[
-2\varepsilon\partial_\theta W_\varepsilon\,\partial_\theta u_\varepsilon
=
2\varepsilon\partial_\theta F_\varepsilon
+
2\varepsilon W_\varepsilon\partial_{\theta\theta}u_\varepsilon .
\]
Consequently, the system can be written as
\begin{equation}\label{eq:WKB_reformulation}
\left\{
\begin{aligned}
& \partial_t u_\varepsilon
- |\partial_\theta u_\varepsilon|^2
- \varepsilon \, \partial_{\theta\theta} u_\varepsilon
= - \mathcal H(\theta,t), \\
& \varepsilon \partial_t W_\varepsilon
- D(\theta)\Delta_x W_\varepsilon
- \varepsilon^2 \partial_{\theta\theta} W_\varepsilon
+2\varepsilon\partial_\theta F_\varepsilon
=
W_\varepsilon
\bigl(
K(x)-\rho_\varepsilon(x,t)
+\mathcal H(\theta,t)
-2\varepsilon\partial_{\theta\theta}u_\varepsilon
\bigr).
\end{aligned}
\right.
\end{equation}
In the continuous theory, \(\mathcal H\) is related to a principal eigenvalue problem in the spatial variable. This structure is the basic reason for introducing the variables \((u_\varepsilon,W_\varepsilon)\) in the first place.

\subsection{Continuous asymptotic structure of the steady state}\label{subsec:continuous_asymptotic_structure}

We next recall the steady-state asymptotic structure that motivates the numerical method. Let \((\bar n_\varepsilon,\bar \rho_\varepsilon)\) be a positive steady state of \eqref{eq:n_intro}. Then
\begin{equation}\label{eq:ss_n}
-D(\theta)\Delta_x \bar n_\varepsilon
-\varepsilon^2\partial_{\theta\theta} \bar n_\varepsilon
=
\bar n_\varepsilon\bigl(K(x)-\bar\rho_\varepsilon(x)\bigr),
\qquad
\bar\rho_\varepsilon(x)=\int_{\mathbb T} \bar n_\varepsilon(x,\theta)\,d\theta.
\end{equation}
As shown in \cite{PerthameSouganidis}, under standard assumptions on \(K\), \(D\), and the domain, the steady-state family admits several properties that are fundamental for the rare-mutation limit.

\begin{proposition}[Continuous asymptotic structure of the steady state]\label{prop:continuous_structure}
Assume that \(K\) is positive and nonconstant, that \(D\) is positive, periodic, and has a unique minimizer \(\theta_m\), and that the steady problem \eqref{eq:ss_n} admits a positive solution. Then the following properties hold at the continuous level.
\begin{enumerate}
\item There exists a constant \(C>0\), independent of \(\varepsilon\), such that
\begin{equation}\label{eq:rho_continuous_bound}
0\le \bar\rho_\varepsilon(x)\le C
\qquad \text{for all } x\in\Omega.
\end{equation}
Moreover, up to subsequences, \(\bar\rho_\varepsilon\) converges uniformly to a limit density \(\rho\).

\item If we write
\begin{equation}\label{eq:u_log_transform}
\bar n_\varepsilon(x,\theta)=\exp\!\left(\frac{\bar u_\varepsilon(x,\theta)}{\varepsilon}\right),
\end{equation}
then \(\bar u_\varepsilon\) is uniformly Lipschitz continuous in \(\theta\), and along subsequences
\begin{equation}\label{eq:u_limit_uniform}
\bar u_\varepsilon \to u
\qquad \text{uniformly in } (x,\theta),
\end{equation}
where the limit \(u\) is independent of \(x\), periodic in \(\theta\), and satisfies
\begin{equation}\label{eq:HJ_limit}
\begin{cases}
|\partial_\theta u(\theta)|^2 = \mathcal H\bigl(\theta,\rho(\cdot)\bigr),\\[1mm]
\displaystyle \max_{\theta\in\mathbb T} u(\theta)=0.
\end{cases}
\end{equation}

\item The effective Hamiltonian \(\mathcal H(\theta,\rho)\) is defined through the principal eigenvalue problem
\begin{equation}\label{eq:eigen_N}
\begin{cases}
-D(\theta)\Delta_x \mathcal N(x,\theta)
= \mathcal N(x,\theta)\bigl(K(x)-\rho(x)\bigr)
+ \mathcal N(x,\theta)\,\mathcal H\bigl(\theta,\rho(\cdot)\bigr), & x\in\Omega,\\
\partial_\nu \mathcal N = 0, & x\in\partial\Omega,
\end{cases}
\end{equation}
with a positive eigenfunction \(\mathcal N\). Moreover, \(\mathcal H\) has the same monotonicity in \(\theta\) as \(D\), and hence
\begin{equation}\label{eq:H_min_zero}
\min_{\theta\in\mathbb T} \mathcal H\bigl(\theta,\rho(\cdot)\bigr)
=
\mathcal H\bigl(\theta_m,\rho(\cdot)\bigr)
=0.
\end{equation}

\item The steady state concentrates on the fittest trait in the limit \(\varepsilon\to0\). More precisely,
\begin{equation}\label{eq:continuous_dirac_limit}
\bar n_\varepsilon \rightharpoonup N_m(x)\,\delta(\theta-\theta_m),
\qquad
\bar\rho_\varepsilon \to N_m(x),
\end{equation}
where \(N_m\) solves the reduced Fisher-type problem
\begin{equation}\label{eq:Nm_problem}
\begin{cases}
-D(\theta_m)\Delta_x N_m = N_m\bigl(K(x)-N_m\bigr), & x\in\Omega,\\
\partial_\nu N_m = 0, & x\in\partial\Omega.
\end{cases}
\end{equation}
\end{enumerate}
\end{proposition}

%\begin{remark}[Interpretation for the numerical method]\label{rem:AP_target_cont}
Proposition~\ref{prop:continuous_structure} provides the continuous target of the numerical scheme. The boundedness of \(\bar\rho_\varepsilon\) keeps the effective Hamiltonian in a bounded regime. The constrained Hamilton--Jacobi equation identifies the dominant trait through the maximizer of \(u\). The concentration result \eqref{eq:continuous_dirac_limit} indicates that an asymptotic-preserving discretization should recover the reduced pair \((N_m,\theta_m)\) without resolving the original density on an extremely fine trait grid. This is the sense in which the WKB reformulation is used in the sequel.
%\end{remark}


%===================================================================================================
%   Detailed scheme
%===================================================================================================

\section{WKB reformulated numerical scheme}

In this section, we develop and analyze a semi-discrete numerical scheme for
the WKB reformulation \eqref{eq:WKB_reformulation}. The scheme is formulated
in terms of the phase variable \(u_\varepsilon\) and the amplitude variable
\(W_\varepsilon\), with the total density \(n_\varepsilon\) reconstructed from
the WKB representation.

\subsection{Semi-discrete numerical discretization}

We work at the semi-discrete level in order to separate the discretization in
\((x,\theta)\) from the additional issue of time stepping. For clarity, we
present the scheme in one spatial dimension with uniform grids in the
\(x\)- and \(\theta\)-directions. This choice is not essential for the WKB
decomposition, but it keeps the notation transparent.

Let \(\{x_j\}_{j=1}^J\) be a uniform grid on \(\Omega=(a,b)\) and let
\(\{\theta_k\}_{k=1}^K\) be a uniform periodic grid on \(\mathbb T=(0,1)\),
with mesh size \(\Delta\theta\). We denote
\[
W_{j,k}^\varepsilon(t)\approx W_\varepsilon(t,x_j,\theta_k),
\qquad
u_k^\varepsilon(t)\approx u_\varepsilon(t,\theta_k).
\]
For notational convenience, we write
\[
D_k=D(\theta_k),
\qquad
K_j=K(x_j).
\]
By \eqref{eq:basic_assumptions}, we have 
\[
    0<D_{\min}\le D_k\le D_{\max}<\infty,
    \qquad
    0\le K_j\le K_{\max}.
\]


For a grid function \(V=\{V_{j,k}\}\), we use the standard second-order
difference operators
\[
\delta_x^2 V_{j,k}
=
\frac{V_{j+1,k}-2V_{j,k}+V_{j-1,k}}{(\Delta x)^2},
\qquad
\delta_\theta^2 V_{j,k}
=
\frac{V_{j,k+1}-2V_{j,k}+V_{j,k-1}}{(\Delta\theta)^2}.
\]
The homogeneous Neumann boundary condition in \(x\) is imposed by the ghost
values
\[
V_{0,k}=V_{1,k},
\qquad
V_{J+1,k}=V_{J,k},
\]
while all indices in the \(\theta\)-direction are understood periodically.
For the phase variable, we also use
\[
D_\theta^-u_k
=
\frac{u_k-u_{k-1}}{\Delta\theta},
\qquad
D_\theta^+u_k
=
\frac{u_{k+1}-u_k}{\Delta\theta}.
\]

The semi-discrete WKB system involves two numerical ingredients that are not
contained in the standard second-order difference operators introduced above.
The first one is the numerical Hamiltonian for the phase equation. The
Hamilton--Jacobi term \(-|\partial_\theta u_\varepsilon|^2\) is approximated by
\[
\widehat H(D_\theta^-u_k^\varepsilon,D_\theta^+u_k^\varepsilon).
\]
We assume that \(\widehat H\) is locally Lipschitz and consistent with the
continuous Hamiltonian, namely
\[
\widehat H(p,p)=-p^2 .
\]
We also assume that it is monotone in the sense that it is nondecreasing with
respect to its first argument and nonincreasing with respect to its second
argument. When \(\widehat H\) is differentiable, this condition reads
\[
\partial_a\widehat H(a,b)\ge0,
\qquad
\partial_b\widehat H(a,b)\le0 .
\]
Typical examples are the Godunov and Lax--Friedrichs numerical Hamiltonians,
as recalled in Appendix~\ref{app:hamiltonian_weno}.

The second ingredient is the conservative discretization of the transport
coupling in the amplitude equation. Recall that the continuous flux is
\[
F_\varepsilon=-W_\varepsilon\partial_\theta u_\varepsilon .
\]
We denote by
\[
\mathcal D_\theta\widehat{\mathcal F}(W^\varepsilon,u^\varepsilon)_{j,k}
\]
a conservative discretization of \(\partial_\theta F_\varepsilon\) at
\((x_j,\theta_k)\). Here \(\widehat{\mathcal F}\) approximates the flux
\(-W_\varepsilon\partial_\theta u_\varepsilon\). We assume that
\(\mathcal D_\theta\widehat{\mathcal F}\) is consistent with
\[
\partial_\theta F_\varepsilon
=
\partial_\theta(-W_\varepsilon\partial_\theta u_\varepsilon)
\]
when evaluated on smooth functions, and that it is conservative in the
periodic trait variable. For the positivity argument below, we further assume
the quasi-positivity property
\begin{equation}\label{ass:H_flux}
W_{j,k}=0,
\qquad
W_{j,\ell}\ge0\ \text{for all }\ell
\quad\Longrightarrow\quad
-\mathcal D_\theta\widehat{\mathcal F}(W,u)_{j,k}\ge0 .
\end{equation}
This property is satisfied by standard monotone finite volume fluxes, such as
the upwind and Lax--Friedrichs fluxes listed in Appendix~\ref{app:fluxes_new}.

With these two ingredients specified, we now define the discrete density and
the semi-discrete WKB system. For notational convenience, set
\begin{equation}\label{def:E}
E_k^\varepsilon(t)
=
\exp\!\left(\frac{u_k^\varepsilon(t)}{\varepsilon}\right).
\end{equation}
The reconstructed nodal density is defined by
\begin{equation}\label{def:n_reconstructed}
n_{j,k}^\varepsilon(t)
=
W_{j,k}^\varepsilon(t)E_k^\varepsilon(t).
\end{equation}


As we shall see later, the evolution of \(u_\varepsilon\) and
\(W_\varepsilon\) is coupled through the spatial marginal density.  In the
semi-discrete scheme, this marginal is computed by the composite quadrature in
the trait direction.  More precisely, we define
\begin{equation}\label{eq:E_rho}
\rho_{j,Q}^\varepsilon(t)
=
m_j^\varepsilon(t)
:=
\Delta\theta
\sum_{k=1}^K
W_{j,k}^\varepsilon(t)E_k^\varepsilon(t),
\qquad
j=1,\dots,J .
\end{equation}
This is the discrete counterpart of
\[
\int_{\mathbb T}
n_\varepsilon(t,x_j,\theta)\,d\theta
=
\int_{\mathbb T}
W_\varepsilon(t,x_j,\theta)
\exp\!\left(\frac{u_\varepsilon(t,\theta)}{\varepsilon}\right)
\,d\theta .
\]



Given
\[
\boldsymbol\rho_Q(t)
=
\bigl(\rho_{1,Q}^\varepsilon(t),\dots,\rho_{J,Q}^\varepsilon(t)\bigr),
\]
the discrete effective Hamiltonian
\(\mathcal H_k(\boldsymbol\rho_Q(t))\) is determined by the discrete
principal eigenvalue problem
\begin{equation}\label{eq:eigen_discrete}
-D_k\delta_x^2\mathcal N_j^{(k)}
=
\mathcal N_j^{(k)}
\bigl(K_j-\rho_{j,Q}^\varepsilon(t)\bigr)
+
\mathcal N_j^{(k)}
\mathcal H_k\bigl(\boldsymbol\rho_Q(t)\bigr),
\qquad
j=1,\dots,J .
\end{equation}
Here \(\mathcal N^{(k)}\) is the corresponding positive eigenvector, equipped
with the same discrete Neumann boundary condition in \(x\). Its normalization
does not affect \(\mathcal H_k\). For definiteness, one may impose
\[
\Delta x\sum_{j=1}^J \mathcal N_j^{(k)}=1 .
\]

With the above notation, the semi-discrete WKB scheme reads
\begin{equation}\label{eq:semi_discrete_WKB_system}
\left\{
\begin{aligned}
&\frac{\mathrm d}{\mathrm dt}u_k^\varepsilon
+
\widehat H
\left(
D_\theta^-u_k^\varepsilon,
D_\theta^+u_k^\varepsilon
\right)
-
\varepsilon\delta_\theta^2u_k^\varepsilon
=
-\mathcal H_k\bigl(\boldsymbol\rho_Q(t)\bigr),
\qquad
k=1,\dots,K,
\\[1mm]
&\varepsilon
\frac{\mathrm d}{\mathrm dt}W_{j,k}^\varepsilon
-
D_k\delta_x^2W_{j,k}^\varepsilon
-
\varepsilon^2\delta_\theta^2W_{j,k}^\varepsilon
+
2\varepsilon
\mathcal D_\theta
\widehat{\mathcal F}(W^\varepsilon,u^\varepsilon)_{j,k}
\\
&\qquad =
W_{j,k}^\varepsilon
\bigl(
K_j-\rho_{j,Q}^\varepsilon(t)
+\mathcal H_k(\boldsymbol\rho_Q(t))
-2\varepsilon\delta_\theta^2u_k^\varepsilon
\bigr),
\qquad
j=1,\dots,J,\quad k=1,\dots,K .
\end{aligned}
\right.
\end{equation}

We next show the equation satisfied by the reconstructed density
\(n_{j,k}^\varepsilon\). Since \(E_k^\varepsilon\) is independent of the
spatial index \(j\),
\[
\delta_x^2 n_{j,k}^\varepsilon
=
E_k^\varepsilon\delta_x^2W_{j,k}^\varepsilon .
\]
Moreover,
\begin{equation}\label{def:dt_n}
\varepsilon\frac{\mathrm d}{\mathrm dt}n_{j,k}^\varepsilon
=
E_k^\varepsilon
\left(
\varepsilon\frac{\mathrm d}{\mathrm dt}W_{j,k}^\varepsilon
+
W_{j,k}^\varepsilon
\frac{\mathrm d}{\mathrm dt}u_k^\varepsilon
\right).
\end{equation}
Multiplying the first equation in \eqref{eq:semi_discrete_WKB_system} by $E_k^\varepsilon W_{j,k}^\varepsilon$ and multiplying the second equation in  \eqref{eq:semi_discrete_WKB_system} by $\varepsilon E_k^\varepsilon$, 
noticing \eqref{def:dt_n},
we can derive that \(n_{j,k}^\varepsilon\) satisfies
\begin{equation}\label{eq:semi_discrete_n}
\varepsilon\frac{\mathrm d}{\mathrm dt} n_{j,k}^\varepsilon
-
D_k\delta_x^2 n_{j,k}^\varepsilon
=
\bigl(K_j-\rho_{j,Q}^\varepsilon\bigr)n_{j,k}^\varepsilon
+
R_{j,k}^\varepsilon ,
\end{equation}
where the remainder has the form 
\begin{align}
R_{j,k}^\varepsilon
&=
\varepsilon^2
E_k^\varepsilon
\delta_\theta^2 W_{j,k}^\varepsilon
-
\varepsilon
W_{j,k}^\varepsilon
E_k^\varepsilon
\delta_\theta^2 u_k^\varepsilon
-
W_{j,k}^\varepsilon
E_k^\varepsilon
\widehat H
\left(
D_\theta^-u_k^\varepsilon,
D_\theta^+u_k^\varepsilon
\right)
-
2\varepsilon
E_k^\varepsilon
\mathcal D_\theta
\widehat{\mathcal F}(W^\varepsilon,u^\varepsilon)_{j,k}.
\label{eq:R_exact}
\end{align}
Thus the reconstructed density satisfies a semi-discrete analogue of the
original density equation, with \(R_{j,k}^\varepsilon\) collecting the
trait-direction discretization effects. When the numerical Hamiltonian and
the conservative flux are evaluated on smooth exact functions, this remainder
is consistent with the trait diffusion contribution
\(\varepsilon^2\partial_{\theta\theta}n^\varepsilon\).





%===================================================================================================
%   Subsection: basic structural property
%===================================================================================================

\subsection{Basic structural property}

The semi-discrete scheme \eqref{eq:semi_discrete_n} can be shown to be positive-preserving 
as long as the initial data is nonnegative. 
The detailed result is summarized as the following lemma. 

\begin{lemma}[Positivity of the amplitude]
\label{lem:positivity_W}
Assume that \(W_{j,k}^\varepsilon(0)\ge0\) for all \(j,k\). Then the
semi-discrete amplitude equation in \eqref{eq:semi_discrete_WKB_system}
preserves nonnegativity, namely
\[
    W_{j,k}^\varepsilon(t)\ge0
    \qquad
    \forall\,j,k,\ t\ge0,
\]
as long as the solution exists. Consequently,
\[
    n_{j,k}^\varepsilon(t)\ge0
    \qquad
    \forall\,j,k,\ t\ge0 .
\]
\end{lemma}

\begin{proof}
It is enough to verify the quasi-positivity of the right-hand side of the
finite-dimensional ODE system for \(W^\varepsilon\). From the amplitude
equation in \eqref{eq:semi_discrete_WKB_system}, we can write
\[
\varepsilon \frac{\mathrm d}{\mathrm dt}W_{j,k}^\varepsilon
=
D_k\delta_x^2W_{j,k}^\varepsilon
+
\varepsilon^2\delta_\theta^2W_{j,k}^\varepsilon
-
2\varepsilon
\mathcal D_\theta\widehat{\mathcal F}(W^\varepsilon,u^\varepsilon)_{j,k}
+
W_{j,k}^\varepsilon B_{j,k}(t),
\]
where
\[
B_{j,k}(t)
=
K_j-\rho_{j,Q}^\varepsilon(t)
+\mathcal H_k(\boldsymbol\rho_Q(t))
-2\varepsilon\delta_\theta^2u_k^\varepsilon(t).
\]
Let \(W^\varepsilon\ge0\) and suppose that
\(W_{j,k}^\varepsilon=0\) for some \((j,k)\). Then, for an interior spatial
point,
\[
\delta_x^2W_{j,k}^\varepsilon
=
\frac{W_{j+1,k}^\varepsilon+W_{j-1,k}^\varepsilon}{(\Delta x)^2}
\ge0, 
\quad 
\delta_\theta^2W_{j,k}^\varepsilon
=
\frac{W_{j,k+1}^\varepsilon+W_{j,k-1}^\varepsilon}{(\Delta\theta)^2}
\ge0 .
\]
The same conclusion holds at the boundary because of the Neumann ghost
values in $x$-direction and the periodicity in \(\theta\)-direction. 
The quasi-positivity assumption on the conservative flux operator implies that 
\[
    -
    \mathcal D_\theta\widehat{\mathcal F}(W^\varepsilon,u^\varepsilon)_{j,k}
    \ge0 .
\]
Combining all the inequalities, noticing the fact that 
\(
    W_{j,k}^\varepsilon B_{j,k}(t)=0 
\)
since we assumed $W_{j,k}^\varepsilon=0$, 
we finally have that, at this specific $(j,k)$, the following inequality holds, 
\[
\left.
\frac{\mathrm d}{\mathrm dt}W_{j,k}^\varepsilon
\right|_{W_{j,k}^\varepsilon=0}
\ge0 .
\]
Thus the vector field points inward on the boundary of the nonnegative cone.
The standard invariance criterion for finite-dimensional ODE systems implies
the claim.
\end{proof}

%\begin{remark}[On the time-dependent stability]
%For each fixed \(\varepsilon>0\), the semi-discrete system is a
%finite-dimensional ODE system under the locally Lipschitz assumptions on the
%discrete operators. Uniform-in-\(\varepsilon\) continuation estimates for the
%time-dependent problem require a separate bootstrap argument involving the
%density, phase, amplitude, and trait-direction remainder. This issue is not
%pursued in the present work.
%\end{remark}


%===================================================================================================
%   Subsection: discrete stationary AP structure
%===================================================================================================

\subsection{Stationary AP structure under a one-sided stability condition}
\label{subsec:stationary_discrete_AP}

We now turn to stationary semi-discrete states. The goal is to identify which
parts of the continuous rare-mutation structure are retained on a fixed grid.
%For the generic split WKB scheme, the only nonstandard point is the
%trait-direction remainder in the reconstructed density equation. We therefore
%state the stationary result under an explicit one-sided stability condition.
%This keeps the AP statement conditional and avoids claiming that the WKB
%remainder is automatically controlled for every trait discretization.
By Lemma~\ref{lem:positivity_W}, stationary states obtained from nonnegative
amplitudes satisfy
\[
    W_{j,k}^\varepsilon\ge0,
    \qquad
    n_{j,k}^\varepsilon
    =W_{j,k}^\varepsilon E_k^\varepsilon\ge0 .
\]
All estimates below are for such nonnegative reconstructed densities.

\paragraph{Weighted density balance.}
Define
\[
    \eta_j^\varepsilon
    =
    \Delta\theta
    \sum_{k=1}^K
    \frac{n_{j,k}^\varepsilon}{D_k} .
\]
Multiplying \eqref{eq:semi_discrete_n} by \(1/D_k\) and summing over the
trait variable gives the stationary weighted balance equation
\begin{equation}\label{eq:weighted_summed_density_stationary}
    -\delta_x^2m_j^\varepsilon
    =
    \eta_j^\varepsilon
    \bigl(K_j-\rho_{j,Q}^\varepsilon\bigr)
    +
    \mathcal R_j^{\varepsilon,D},
\end{equation}
where
\[
    m_j^\varepsilon
    =
    \Delta\theta\sum_{k=1}^K n_{j,k}^\varepsilon,
    \qquad
    \mathcal R_j^{\varepsilon,D}
    =
    \Delta\theta
    \sum_{k=1}^K
    \frac{R_{j,k}^\varepsilon}{D_k} .
\]
The term \(\mathcal R_j^{\varepsilon,D}\) is the discrete analogue of the
continuous weighted mutation contribution
\[
    \varepsilon^2
    \int_{\mathbb T}
    \frac1{D(\theta)}
    \partial_{\theta\theta}n^\varepsilon
    \,d\theta
    =
    \varepsilon^2
    \int_{\mathbb T}
    n^\varepsilon
    \partial_{\theta\theta}\left(\frac1{D(\theta)}\right)
    \,d\theta .
\]
At the continuous level this is bounded above by
\(C\varepsilon^2\rho^\varepsilon\). For the split WKB discretization, however,
the discrete operators act on the reconstructed density
\(n_{j,k}^\varepsilon=W_{j,k}^\varepsilon E_k^\varepsilon\), and an exact
chain rule is not available. We isolate the needed property as follows.

\begin{assumption}[One-sided weighted remainder stability]
\label{ass:remainder_stability}
There exist constants \(C_R>0\) and \(r_h\ge0\), independent of
\(\varepsilon\), such that every stationary state under consideration
satisfies
\begin{equation}\label{eq:remainder_CR_bound}
    \bigl(\mathcal R_j^{\varepsilon,D}\bigr)_+
    \le
    C_R m_j^\varepsilon+r_h,
    \qquad j=1,\ldots,J .
\end{equation}
\end{assumption}

This condition only controls the positive part of the weighted defect; a
negative contribution is favorable in the maximum-principle estimate. It is a
structural stability requirement on the trait discretization, not a consequence
of the preceding consistency statement. For a direct density-level conservative
discretization, and for the density-compatible WKB variant
\eqref{eq:density_compatible_W}, the same type of estimate follows by discrete
summation by parts. For the generic split WKB scheme, it should be checked
analytically under additional structure or monitored numerically.

We also record the mild stability condition needed for the density
reconstruction in the nonlocal reaction term.

\begin{assumption}[Reconstruction stability]
\label{ass:reconstruction_stability}
There exist constants \(c_Q,C_Q>0\) and \(e_Q\ge0\), independent of
\(\varepsilon\), such that
\begin{equation}\label{eq:quadrature_coercivity_lower}
    \rho_{j,Q}^\varepsilon
    \ge
    c_Qm_j^\varepsilon-e_Q,
    \qquad j=1,\ldots,J,
\end{equation}
\begin{equation}\label{eq:quadrature_coercivity_upper}
    \rho_{j,Q}^\varepsilon
    \le
    C_Qm_j^\varepsilon+e_Q,
    \qquad j=1,\ldots,J .
\end{equation}
\end{assumption}
For the direct composite quadrature \(\rho_{j,Q}^\varepsilon=m_j^\varepsilon\),
Assumption~\ref{ass:reconstruction_stability} holds with
\(c_Q=C_Q=1\) and \(e_Q=0\).

\begin{proposition}[Conditional density bound for stationary semi-discrete states]
\label{prop:uniform_bound_stationary_mass}
Let \((u^\varepsilon,W^\varepsilon)\) be a stationary semi-discrete state with
nonnegative reconstructed density. Assume that
Assumptions~\ref{ass:remainder_stability} and
\ref{ass:reconstruction_stability} hold. Then there exist constants
\(C_m,C_\rho>0\), independent of \(\varepsilon\), such that
\[
    0\le m_j^\varepsilon\le C_m,
    \qquad
    0\le \rho_{j,Q}^\varepsilon\le C_\rho,
    \qquad j=1,\ldots,J .
\]
\end{proposition}

\begin{proof}
Since \(n_{j,k}^\varepsilon\ge0\),
\[
    \frac1{D_{\max}}m_j^\varepsilon
    \le
    \eta_j^\varepsilon
    \le
    \frac1{D_{\min}}m_j^\varepsilon .
\]
Using \eqref{eq:weighted_summed_density_stationary} and
\eqref{eq:remainder_CR_bound}, we obtain
\[
    -\delta_x^2m_j^\varepsilon
    +
    \eta_j^\varepsilon\rho_{j,Q}^\varepsilon
    \le
    K_{\max}\eta_j^\varepsilon
    +
    C_Rm_j^\varepsilon
    +
    r_h .
\]
Let \(j_*\) be a point where
\(M^\varepsilon:=\max_jm_j^\varepsilon=m_{j_*}^\varepsilon\). Since
\(\delta_x^2m_{j_*}^\varepsilon\le0\),
\[
    \eta_{j_*}^\varepsilon\rho_{j_*,Q}^\varepsilon
    \le
    K_{\max}\eta_{j_*}^\varepsilon
    +
    C_RM^\varepsilon+r_h .
\]
If \(M^\varepsilon\le e_Q/c_Q\), the bound is immediate. Otherwise,
\(\rho_{j_*,Q}^\varepsilon\ge c_QM^\varepsilon-e_Q>0\), and hence
\[
    \frac{M^\varepsilon}{D_{\max}}
    \bigl(c_QM^\varepsilon-e_Q\bigr)
    \le
    \left(\frac{K_{\max}}{D_{\min}}+C_R\right)M^\varepsilon+r_h .
\]
This quadratic inequality gives a bound for \(M^\varepsilon\) independent of
\(\varepsilon\). The upper estimate for \(\rho_{j,Q}^\varepsilon\) then follows
from \eqref{eq:quadrature_coercivity_upper}.
\end{proof}

\paragraph{Stationary phase estimate.}
The density bound controls the coefficients in the discrete eigenvalue problem
and hence the stationary phase equation.

\begin{proposition}[Stationary phase estimate on a fixed trait grid]
\label{prop:stationary_phase_estimate}
Assume that
\[
    0\le \rho_{j,Q}^\varepsilon\le \overline\rho,
    \qquad j=1,\ldots,J,
\]
where \(\overline\rho\) is independent of \(\varepsilon\). Then the discrete
effective Hamiltonian defined by \eqref{eq:eigen_discrete} satisfies
\[
    |\mathcal H_k(\boldsymbol\rho_Q)|
    \le C_{\mathcal H},
    \qquad k=1,\ldots,K,
\]
where \(C_{\mathcal H}\) is independent of \(\varepsilon\).

Assume further that the stationary phase equation
\begin{equation}\label{eq:stationary_phase}
    \widehat H_G
    \left(
        D_\theta^-u_k^\varepsilon,
        D_\theta^+u_k^\varepsilon
    \right)
    -
    \varepsilon\delta_\theta^2u_k^\varepsilon
    =
    -\mathcal H_k(\boldsymbol\rho_Q)
\end{equation}
holds with the Godunov numerical Hamiltonian
\[
    \widehat H_G(p^-,p^+)
    =
    -\bigl((p^-)_{-}\bigr)^2
    -
    \bigl((p^+)_{+}\bigr)^2 .
\]
Then
\[
    \Delta\theta
    \sum_{k=1}^K
    |D_\theta^+u_k^\varepsilon|^2
    \le
    C_{\mathcal H} .
\]
Consequently,
\[
    \max_k |D_\theta^+u_k^\varepsilon|
    \le
    \frac{\sqrt{C_{\mathcal H}}}{\sqrt{\Delta\theta}} .
\]
If the phase is normalized by
\[
    \max_k u_k^\varepsilon=0,
\]
then
\[
    \max_k |u_k^\varepsilon|
    \le
    \sqrt{C_{\mathcal H}} .
\]
All constants are independent of \(\varepsilon\), although they may depend on
the fixed trait mesh.
\end{proposition}

\begin{proof}
The discrete eigenvalue problem is associated with the symmetric operator
\[
    -D_k\delta_x^2-(K_j-\rho_{j,Q}^\varepsilon).
\]
Its principal eigenvalue has the Rayleigh representation
\[
    \mathcal H_k(\boldsymbol\rho_Q)
    =
    \min_{\varphi\ne0}
    \frac{
        D_k\Delta x\sum_j |\delta_x^+\varphi_j|^2
        -
        \Delta x\sum_j
        (K_j-\rho_{j,Q}^\varepsilon)\varphi_j^2
    }{
        \Delta x\sum_j\varphi_j^2
    } .
\]
Since \(0\le K_j\le K_{\max}\) and
\(0\le\rho_{j,Q}^\varepsilon\le\overline\rho\), we have
\[
    \mathcal H_k(\boldsymbol\rho_Q)\ge -K_{\max} .
\]
Testing the quotient with a constant vector gives
\(\mathcal H_k(\boldsymbol\rho_Q)\le\overline\rho\). Thus
\[
    |\mathcal H_k(\boldsymbol\rho_Q)|
    \le
    \max\{K_{\max},\overline\rho\}
    =:C_{\mathcal H}.
\]

Set \(q_k^\varepsilon=D_\theta^+u_k^\varepsilon\). Then
\(D_\theta^-u_k^\varepsilon=q_{k-1}^\varepsilon\), and
\[
    -\widehat H_G
    \left(
        D_\theta^-u_k^\varepsilon,
        D_\theta^+u_k^\varepsilon
    \right)
    =
    \bigl((q_{k-1}^\varepsilon)_{-}\bigr)^2
    +
    \bigl((q_k^\varepsilon)_{+}\bigr)^2 .
\]
Summing \eqref{eq:stationary_phase} over the periodic trait grid eliminates
\(\delta_\theta^2u^\varepsilon\). Therefore
\[
\begin{aligned}
    \Delta\theta\sum_{k=1}^K |q_k^\varepsilon|^2
    &=
    -\Delta\theta\sum_{k=1}^K
    \widehat H_G
    \left(
        D_\theta^-u_k^\varepsilon,
        D_\theta^+u_k^\varepsilon
    \right)  \\
    &=
    \Delta\theta\sum_{k=1}^K
    \mathcal H_k(\boldsymbol\rho_Q)
    \le C_{\mathcal H} .
\end{aligned}
\]
The pointwise slope bound follows from the fixed-grid inequality
\(|q_k|^2\le(\Delta\theta)^{-1}\Delta\theta\sum_\ell |q_\ell|^2\). If
\(\max_k u_k^\varepsilon=0\), connecting any node to a maximizer along the
periodic grid gives
\[
    |u_k^\varepsilon|
    \le
    \Delta\theta\sum_{\ell=1}^K |D_\theta^+u_\ell^\varepsilon|
    \le
    \left(
    \Delta\theta\sum_{\ell=1}^K |D_\theta^+u_\ell^\varepsilon|^2
    \right)^{1/2}.
\]
This proves the claim.
\end{proof}

\paragraph{Formal AP trait-selection limit.}
We finally pass formally to the fixed-grid rare-mutation limit. The statement
is deliberately fixed-grid and conditional: it identifies the limiting
selection mechanism once the density and phase estimates above are available.

\begin{proposition}[Formal AP trait-selection limit]
\label{prop:formal_AP_trait_selection}
Assume that the stationary semi-discrete states satisfy
Assumptions~\ref{ass:remainder_stability} and
\ref{ass:reconstruction_stability}, and that the phase is normalized by
\[
    \max_k u_k^\varepsilon=0 .
\]
Assume also that the stationary phase estimate of
Proposition~\ref{prop:stationary_phase_estimate} applies. Then, up to a
subsequence on the fixed grids,
\[
    \boldsymbol\rho_Q^\varepsilon\to \boldsymbol\rho^0,
    \qquad
    u^\varepsilon\to u^0 .
\]
Moreover, the limiting phase satisfies the discrete constrained
Hamilton--Jacobi system
\[
    \widehat H
    \left(
        D_\theta^-u_k^0,
        D_\theta^+u_k^0
    \right)
    =
    -\mathcal H_k(\boldsymbol\rho^0),
    \qquad
    \max_k u_k^0=0 .
\]
For the Godunov Hamiltonian \(\widehat H_G\), every maximizer \(k_*\) of
\(u^0\) satisfies
\[
    \mathcal H_{k_*}(\boldsymbol\rho^0)=0 .
\]
If \(\mathcal H_k(\boldsymbol\rho^0)\) has a unique zero at \(k_m\), then
\(k_*=k_m\). Hence the selected grid trait is identified by
\[
    \theta_{k_m}\in \arg\max_k u_k^0 .
\]
\end{proposition}

\begin{proof}
Proposition~\ref{prop:uniform_bound_stationary_mass} and
\eqref{eq:quadrature_coercivity_upper} give a uniform bound on
\(\boldsymbol\rho_Q^\varepsilon\). Since the grids are fixed, a subsequence
converges to some \(\boldsymbol\rho^0\). The normalized phase is also uniformly
bounded on the fixed trait grid by
Proposition~\ref{prop:stationary_phase_estimate}; hence, up to a further
subsequence, \(u^\varepsilon\to u^0\).

Passing to the limit in the stationary phase equation gives the discrete
constrained Hamilton--Jacobi system, because
\(\varepsilon\delta_\theta^2u_k^\varepsilon\to0\) on the fixed grid. For the
Godunov Hamiltonian, \(\widehat H_G\le0\), and therefore
\(
\mathcal H_k(\boldsymbol\rho^0)=-\widehat H_G(D_\theta^-u_k^0,D_\theta^+u_k^0)
\ge0
\).
At a maximizer \(k_*\) of \(u^0\),
\[
    D_\theta^-u_{k_*}^0\ge0,
    \qquad
    D_\theta^+u_{k_*}^0\le0,
\]
so \(\widehat H_G(D_\theta^-u_{k_*}^0,D_\theta^+u_{k_*}^0)=0\). Hence
\(\mathcal H_{k_*}(\boldsymbol\rho^0)=0\). If this zero is unique, the maximizer
must be \(k_m\).
\end{proof}




\begin{remark}
The above result describes the fixed-grid AP structure of the stationary
selection state.  In the computations, this stationary state is reached through
the time-dependent WKB evolution.  Thus the evolutionary scheme should be viewed
as a numerical relaxation path toward the discrete eigenvalue and constrained
Hamilton--Jacobi structure characterized above.
This distinction is useful for interpreting the fully discrete method.  The
semi-discrete analysis identifies the limiting stationary structure, while the
fully discrete algorithm evolves \(u\) and \(W\) toward this structure. 
\end{remark}








\section{Fully discrete scheme and reconstruction techniques}
\label{subsec:fully_discrete_reconstruction}

In this section, we describe the fully discrete WKB method and the
reconstruction procedures used to recover the trait concentration profiles.  
The unknowns advanced in time are the WKB variables \(u\) and \(W\).  
The density, the spatial marginal, and the trait marginal are then recovered from these
variables through a reconstruction layer. 
 This layer is part of the numerical
method rather than a visualization step.  In the small mutation regime, direct
evaluation of \(W\exp(u/\varepsilon)\) may suffer from overflow, underflow, and
exponential amplification of small phase errors.  The reconstruction must
therefore be stable, mass-consistent, and computationally efficient.

\subsection{Fully discrete scheme}

Let \(t_n=n\tau\).  Assume that \(u^n\) and \(W^n\) are known at time \(t_n\).
The reconstructed density and spatial marginal are denoted by
\[
    n^n={\mathcal R}_n[W^n,u^n],
    \qquad
    \rho^n={\mathcal R}_\rho[W^n,u^n].
\]
After \(\rho^n\) is obtained, the discrete effective Hamiltonian
\(H^n=\mathcal H[\rho^n]\) is computed from the corresponding discrete
principal eigenvalue problem.

The phase variable is advanced by
\begin{equation}
    \frac{\widetilde u_k^{n+1}-u_k^n}{\tau}
    +
    \widehat H(D_\theta^-u_k^n,D_\theta^+u_k^n)
    -
    \varepsilon\delta_\theta^2\widetilde u_k^{n+1}
    =
    -H_k^n .
    \label{eq:fully-discrete-u}
\end{equation}
Here \(\widehat H\) is the monotone numerical Hamiltonian used in the
semi-discrete scheme.  The diffusion term in the trait direction is treated
implicitly in order to avoid a severe parabolic time-step restriction.

With \(\widetilde u^{n+1}\) fixed, the amplitude is first advanced by a
provisional update.  A typical split update takes the form
\begin{equation}
\begin{aligned}
    &\varepsilon\frac{\widetilde W_{j,k}^{n+1}-W_{j,k}^n}{\tau}
    -D_k\delta_x^2\widetilde W_{j,k}^{n+1}
    -\varepsilon^2\delta_\theta^2\widetilde W_{j,k}^{n+1}
    +2\varepsilon D_\theta\widehat F(\widetilde W^{n+1},
    \widetilde u^{n+1})_{j,k}
    \\
    &\qquad
    =
    \widetilde W_{j,k}^{n+1}
    \left(
    K_j-\rho_j^n+H_k^n-2\varepsilon\delta_\theta^2
    \widetilde u_k^{n+1}
    \right).
\end{aligned}
\label{eq:fully-discrete-W-provisional}
\end{equation}
The numerical flux \(\widehat F\) is chosen consistently with the conservative
trait-transport term in the semi-discrete formulation.  Density-compatible and
Patankar-type variants of the amplitude update may also be used as fully
discrete stabilization devices.

After the provisional update, we apply a scalar mass-balance projection.  This
step is motivated by the density-level identity
\[
    \varepsilon \frac{dM}{dt}
    =
    R(\rho),
    \qquad
    M(t)=\int_\Omega \rho(x,t)\,dx,
    \qquad
    R(\rho)=\int_\Omega \rho(x,t)(K(x)-\rho(x,t))\,dx .
\]
Let \(\widetilde\rho^{n+1}\) be the spatial marginal reconstructed from
\((\widetilde W^{n+1},\widetilde u^{n+1})\).  We determine a positive scalar
\(\lambda^{n+1}\) from a discrete mass-balance relation and set
\[
    W^{n+1}=\lambda^{n+1}\widetilde W^{n+1},
    \qquad
    u^{n+1}=\widetilde u^{n+1}.
\]
Thus the density-level quantities reconstructed from \((W^{n+1},u^{n+1})\)
are rescaled by the same factor, while the phase profile is kept fixed.  The
projection therefore does not move the selected trait location, but it improves
the consistency of the reconstructed mass and stabilizes the peak height of the
trait marginal.  The algebraic form of the scalar projection is given in
Appendix~\ref{app:mass_correction}.

Finally, a gauge normalization may be applied to \(u^{n+1}\) and \(W^{n+1}\).
Since the WKB representation is invariant under the transformation
\[
    u\mapsto u-c,
    \qquad
    W\mapsto W\exp(c/\varepsilon),
\]
this normalization improves numerical conditioning without changing the
reconstructed density.  The stopping criterion is based on gauge-invariant
residuals involving the normalized phase and the reconstructed density.

\subsection{Reconstruction techniques}

The reconstruction layer has three main purposes.  First, it reconstructs the
solver-level spatial marginal \(\rho(x)\), which enters the nonlocal feedback
and the effective Hamiltonian.  Second, it reconstructs the trait marginal
\(P(\theta)\), from which the selected trait, peak height, and peak width are
extracted.  Third, it enables a dual trait-grid implementation that resolves
the one-dimensional phase more finely than the space-dependent amplitude, and
thereby reduces the dominant computational cost.

\paragraph{Reconstruction of the spatial marginal.}

For each spatial grid point \(x_j\), define
\[
    \ell_{j,k}^n
    =
    \log\bigl(\max\{W_{j,k}^n,W_{\rm floor}\}\bigr)
    +
    \frac{u_k^n}{\varepsilon}.
\]
A stable log-sum-exp reconstruction of the spatial marginal is then given by
\begin{equation}
    \widetilde\rho_j^n
    =
    \exp(m_j^n)
    \sum_k \omega_k^\theta \exp(\ell_{j,k}^n-m_j^n),
    \qquad
    m_j^n=\max_k \ell_{j,k}^n .
    \label{eq:rho-logsumexp}
\end{equation}
This formula is algebraically equivalent to the direct quadrature of
\(W\exp(u/\varepsilon)\), but it avoids overflow and underflow caused by large
values of \(u/\varepsilon\).

When the trait profile is highly concentrated, the quadrature in
\eqref{eq:rho-logsumexp} may become sensitive to the trait mesh.  In this case,
one may use a local Laplace reconstruction near the maximizer of the phase.
This reconstruction uses the local curvature of \(u\) to approximate the
concentrated contribution of the exponential factor.  It is useful when the
physical width of the concentration layer is smaller than the trait mesh size.

\paragraph{Reconstruction of the trait marginal.}

The trait marginal is reconstructed from the corrected WKB variables,
\[
    P_k^n
    =
    \int_\Omega W^n(x,\theta_k)
    \exp(u_k^n/\varepsilon)\,dx .
\]
Set
\[
    A_k^n=\int_\Omega W^n(x,\theta_k)\,dx,
    \qquad
    L_k^n
    =
    \log P_k^n
    =
    \frac{u_k^n}{\varepsilon}
    +
    \log\bigl(\max\{A_k^n,A_{\rm floor}\}\bigr).
\]
The reconstruction is performed on \(L_k^n=\log P_k^n\), rather than on
\(P_k^n\) itself.  This is important because \(P_k^n\) is exponentially
sensitive to errors in \(u_k^n\), whereas \(L_k^n\) remains a smooth quantity on
the WKB scale.

The selected trait is obtained by a two-step local reconstruction.  First,
\(L_k^n\) is reconstructed in the trait direction by a WENO-type interpolation.
Second, a local parabolic refinement is applied near the reconstructed maximum.
This gives the peak location \(\theta_m^n\).  Since the scalar mass-balance
projection changes only the density scale, it affects the reconstructed peak
height but not the phase-based peak location.  The same local quadratic
approximation also provides stable estimates of the peak height and the peak
width.  In this way, the narrow peak of \(P(\theta)\) can be recovered even when
the original trait grid is relatively coarse.

\paragraph{Dual trait-grid implementation.}

Another important component of the fully discrete method is the dual trait-grid
strategy.  The phase \(u\) and the effective Hamiltonian are computed on a
relatively fine trait grid, whereas the space-dependent amplitude \(W\) may be
evolved on a coarser trait grid.  This separation is motivated by the different
dimensions and roles of the two WKB variables.  The phase \(u\) depends only on
the trait variable and controls the location and sharpness of the
concentration.  It is therefore inexpensive to resolve \(u\) on a finer trait
grid.  By contrast, \(W\) depends on both space and trait.  In \(d\) spatial
dimensions, the storage and evolution of \(W\) scale like
\[
    O(N_x^d N_\theta^W),
\]
up to the cost of the spatial solvers, while the phase variable only carries a
one-dimensional trait cost.  Hence reducing \(N_\theta^W\) can lead to a
substantial saving in memory and CPU time, especially when the spatial dimension
or the spatial resolution is large.

In the dual-grid implementation, the fine-grid phase is used to resolve the
trait-scale information and to locate the peak, while the amplitude is stored
and evolved on the coarser trait grid needed for the spatial profile.  Transfer
between the two trait grids is performed within the reconstruction layer.  This
design preserves the main advantage of the WKB formulation: the concentrated
trait profile is resolved through the low-dimensional phase rather than by
evolving the full space-trait density on a fine trait mesh.

In summary, the fully discrete method consists of a time-marching step for the
WKB variables, a scalar mass-balance projection for density-level consistency,
and a reconstruction layer for \(\rho\), \(P(\theta)\), and the local peak
structure.  The log-stabilized reconstruction improves robustness in the small
mutation regime, the local peak reconstruction gives accurate selected-trait
diagnostics, and the dual-grid strategy reduces the dominant cost associated
with the space-dependent amplitude.



%===================================================================================================
%	Numerical experiments
%===================================================================================================


\section{Numerical experiments}\label{sec:numerics}

This section describes the numerical experiments used to test the WKB formulation.  The tests are organized to match the main assumptions and claims of the method: the structural assumptions used in the stationary AP argument, the emergence of concentration in the trait direction, the accuracy of the reconstructed density when the grid is sufficiently fine, the advantage of the WKB phase on coarse trait grids, standard accuracy checks, and a two-dimensional spatial experiment.

\subsection{Common computational setting}\label{subsec:numerical_common}

Unless otherwise stated, the one-dimensional tests use
\[
    \Omega=(0,1),\qquad \mathbb T=[0,1),
\]
with homogeneous Neumann boundary conditions in the spatial variable and periodic boundary conditions in the trait variable.  The baseline coefficients used in the current one-dimensional code are
\begin{equation}\label{eq:num_K_baseline_new}
    K(x)=K_0-K_1\cos(2\pi x),
    \qquad K_0=1,
    \qquad K_1=0.5,
\end{equation}
\begin{equation}\label{eq:num_D_baseline_new}
    D(\theta)=D_{\min}+D_1\bigl(1-\cos(2\pi(\theta-\theta_m))\bigr),
    \qquad D_{\min}=0.2,
    \qquad D_1=0.4,
    \qquad \theta_m=0.35.
\end{equation}
Thus the dispersal rate has a unique minimum at \(\theta_m\), which is the expected selected trait in the rare-mutation limit.  The run scripts specify the initial phase for each test.  The two representative choices are a flat phase and a peaked phase,
\begin{equation}\label{eq:num_initial_data_new}
    u_k^0=0
    \qquad\hbox{or}\qquad
    u_k^0=-\alpha_0\bigl(1-\cos(2\pi(\theta_k-\theta_0))\bigr),
    \quad \theta_0=0.70,
\end{equation}
with \(\alpha_0=0.05\) in the non-well-prepared tests and
\(\alpha_0=O(\varepsilon)\) in some small-\(\varepsilon\) refinement tests.  The amplitude is usually initialized by
\begin{equation}\label{eq:num_initial_W_new}
    W_{j,k}^0=1,
    \qquad
    n_{j,k}^0=W_{j,k}^0\exp(u_k^0/\varepsilon).
\end{equation}
When the peaked phase is used, its initial maximizer \(\theta_0\) is deliberately placed away from \(\theta_m\), so the computation tests selection rather than preservation of a well-prepared peak.

For the WKB computation we use the phase equation in \eqref{eq:semi_discrete_WKB_system} together with the selected amplitude update.  The density in the nonlocal reaction term is reconstructed from \((W^n,u^n)\).  The log-stabilized direct quadrature is
\begin{equation}\label{eq:log_stable_rho_new}
    \rho_{j,Q}^n
    =\exp(s_j^n)\Delta\theta
    \sum_{k=1}^{K} \exp\left(
    \log\bigl(\max\{W_{j,k}^n,W_{\rm floor}\}\bigr)
    +\frac{u_k^n}{\varepsilon}-s_j^n\right),
\end{equation}
where
\[
    s_j^n=\max_k\left
    \{\log\bigl(\max\{W_{j,k}^n,W_{\rm floor}\}\bigr)
    +\frac{u_k^n}{\varepsilon}\right\}.
\]
The code also uses the hybrid Laplace reconstruction described in
Subsection~\ref{subsec:fully_discrete_time} and Appendix~\ref{app:reconstruction}; in particular, the small-\(\varepsilon\) runs can use Laplace reconstruction for \(\rho\), while the trait marginal \(P(\theta)\) is reconstructed in logarithmic form for peak localization and width estimation.  After each time step, we use the gauge normalization
\begin{equation}\label{eq:gauge_new}
    c^n=\max_k u_k^n,
    \qquad
    u_k^n\leftarrow u_k^n-c^n,
    \qquad
    W_{j,k}^n\leftarrow W_{j,k}^n\exp(c^n/\varepsilon),
\end{equation}
which preserves the reconstructed density \(W_{j,k}^n\exp(u_k^n/\varepsilon)\). {\color{red} The fully discrete reconstruction pipeline, residuals, time-step restrictions, mass correction, and dual trait-grid option used in the code are summarized in Subsection~\ref{subsec:fully_discrete_time}.}

The raw trait marginal, selected traits, and concentration width are computed as
\begin{equation}\label{eq:diagnostics_trait_new}
    P_k=\int_\Omega n(x,\theta_k)\,dx,
    \qquad
    \theta_{\rm dir}=\theta_{\arg\max_k P_k},
    \qquad
    \theta_{\rm WKB}=\theta_{\arg\max_k u_k},
\end{equation}
while the postprocessed small-\(\varepsilon\) marginal uses the logP-WENO--mass--quadratic reconstruction described in \eqref{eq:logP_reconstruction_new}--\eqref{eq:logP_quadratic_width_new}.  In that case the reported peak location is the reconstructed maximizer of \(\log P\), with optional parabolic subgrid refinement, rather than the maximizer of a linearly interpolated marginal.
\begin{equation}\label{eq:periodic_distance_new}
    d_{\mathbb T}(a,b)=\min_{q\in\mathbb Z}|a-b+q|,
\end{equation}
\begin{equation}\label{eq:sigma_theta_new}
    \sigma_\theta=
    \left(
    \frac{\sum_{k=1}^{K}d_{\mathbb T}(\theta_k,\theta_{\rm WKB})^2P_k}
         {\sum_{k=1}^{K}P_k}
    \right)^{1/2}.
\end{equation}
For the weighted remainder assumption, we monitor
\begin{equation}\label{eq:gamma_R_new}
    \Gamma_R^\varepsilon
    =
    \max_{1\le j\le J}
    \frac{(R_j^{\varepsilon,D})_+}{m_j^\varepsilon+10^{-14}},
    \qquad
    m_j^\varepsilon=\Delta\theta\sum_{k=1}^{K}n_{j,k}^\varepsilon,
    \qquad
    R_j^{\varepsilon,D}=\Delta\theta\sum_{k=1}^{K}\frac{R_{j,k}^\varepsilon}{D_k}.
\end{equation}
A bounded value of \(\Gamma_R^\varepsilon\) as \(\varepsilon\) decreases is not a proof of Assumption~\ref{ass:remainder_stability}, but it is a useful numerical diagnostic for the positive part of the weighted defect. {\color{red} In addition to \(\Gamma_R^\varepsilon\), the fully discrete diagnostics described in Subsection~\ref{subsec:fully_discrete_time} are reported. In particular, the phase residual, amplitude residual, density residual, and the selected trait are monitored separately, since the WKB phase may converge to the selected trait before the full amplitude has reached a numerical steady state.} For the direct density solver, the analogous residual is computed with \(n^{n+1}-n^n\).

\subsection{Test 1: numerical verification of the structural assumptions}\label{subsec:test_assumptions_new}

The first test checks the structural quantities that appear in the stationary AP argument.  We use
\[
    J=200,
    \qquad
    K=64,
    \qquad
    \varepsilon\in\{5\times10^{-2},2\times10^{-2},10^{-2},5\times10^{-3}\}.
\]
Both the split WKB update and the density-compatible update are run.  The following quantities are reported:
\begin{itemize}
\item \(\min_{j,k} W_{j,k}\), to monitor positivity of the amplitude;
\item \(\max_j\rho_{j,Q}\), to monitor the density bound;
\item \(\Gamma_R^\varepsilon\), to test the one-sided weighted remainder condition;
\item \(\min_k H_k(\rho_Q)\) and the location of this minimum, to check consistency with the selected dispersal minimizer;
\item \(d_{\mathbb T}(\theta_{\rm WKB},\theta_m)\), to check trait selection.
\end{itemize}
The density-compatible update is expected to keep \(\Gamma_R^\varepsilon\) controlled by the discrete second difference of \(1/D_k\), as explained in Appendix~\ref{app:remainder}.  For the generic split update, this table is a diagnostic for the conditional assumption.

\begin{table}[ht]
\centering
\caption{Test 1. Numerical diagnostics for the structural assumptions.}
\begin{tabular}{ccccccc}
\toprule
variant & $\varepsilon$ & $\min W$ & $\max\rho_Q$ & $\Gamma_R^\varepsilon$ & $\theta_{\rm WKB}$ & $d_{\mathbb T}(\theta_{\rm WKB},\theta_m)$\\
\midrule
split & $5\times10^{-2}$ & -- & -- & -- & -- & --\\
split & $2\times10^{-2}$ & -- & -- & -- & -- & --\\
split & $10^{-2}$ & -- & -- & -- & -- & --\\
split & $5\times10^{-3}$ & -- & -- & -- & -- & --\\
density-compatible & $5\times10^{-2}$ & -- & -- & -- & -- & --\\
density-compatible & $2\times10^{-2}$ & -- & -- & -- & -- & --\\
density-compatible & $10^{-2}$ & -- & -- & -- & -- & --\\
density-compatible & $5\times10^{-3}$ & -- & -- & -- & -- & --\\
\bottomrule
\end{tabular}
\end{table}

\subsection{Test 2: long-time concentration in the trait direction}\label{subsec:test_long_time_concentration_new}

This test checks whether the solution actually concentrates in the trait direction after sufficiently large time.  We use
\[
    J=200,
    \qquad K=128,
    \qquad \varepsilon\in\{5\times10^{-2},2\times10^{-2},10^{-2}\},
\]
and store snapshots at several times, for example
\[
    t\in\{1,5,10,20,40,80\}.
\]
For each snapshot, we plot the trait marginal \(P_k(t)\), the phase \(u_k(t)\), and record \(\theta_{\rm WKB}(t)\), \(\theta_{\rm dir}(t)\), and \(\sigma_\theta(t)\).  The expected behavior is that the maximizer of \(u\) moves toward \(\theta_m\), while the concentration width \(\sigma_\theta\) decreases with time and with \(\varepsilon\).

\begin{table}[ht]
\centering
\caption{Test 2. Long-time concentration diagnostics.}
\begin{tabular}{cccccc}
\toprule
$\varepsilon$ & $t$ & $\theta_{\rm WKB}(t)$ & $\theta_{\rm dir}(t)$ & $\sigma_\theta(t)$ & $\min_k H_k(\rho_Q(t))$\\
\midrule
$5\times10^{-2}$ & -- & -- & -- & -- & --\\
$2\times10^{-2}$ & -- & -- & -- & -- & --\\
$10^{-2}$ & -- & -- & -- & -- & --\\
\bottomrule
\end{tabular}
\end{table}


\subsection{Test 3: small-\(\varepsilon\) trait-marginal reconstruction and dual-grid WKB advantage}\label{subsec:test_fine_grid_n_new}

This test focuses on the quantity that becomes most singular when
\(\varepsilon\) is small: the trait marginal \(P(\theta)\).  The current code
uses a very small mutation parameter, for example
\[
    \varepsilon=10^{-4},
    \qquad J=40,
    \qquad T=2,
\]
and compares several WKB configurations with coarse amplitude trait grids.  In
the dual-grid runs, the amplitude grid has \(K_W=8\) or \(16\) points, while the
phase grid has \(K_u=32\) or \(64\) points.  Thus the phase peak and curvature
are resolved on a finer one-dimensional grid, but the space-dependent amplitude
is still advanced on a much coarser trait grid.

For each saved snapshot we compare three reconstructions of \(P(\theta)\):
\[
\begin{array}{ll}
{\rm raw:} & \hbox{compute the nodal marginal and do no peak reconstruction},\\[1mm]
{\rm basic:} & \hbox{interpolate \(\log P\) by a standard periodic interpolant},\\[1mm]
{\rm full:} & \hbox{use WENO peak localization in \(\log P\), mass normalization,}
              \\
              & \hbox{and the phase-amplitude quadratic reconstruction
              \eqref{eq:logP_quadratic_width_new}.}
\end{array}
\]
The purpose is to separate two errors that are otherwise mixed together.  The
first is the error in the WKB variables \((W,u)\); the second is the error made
when a very sharp marginal is reconstructed from those variables.  Directly
interpolating \(P\) is not reliable in this regime, because \(P\) may be
localized between two coarse trait nodes and its height is exponentially
sensitive to the phase.  Reconstructing \(\log P\), locating the peak by WENO,
and then imposing the corrected mass stabilizes the peak position and height.
The local quadratic fit supplies the peak-width estimate \(\sigma_P\).

\begin{table}[ht]
\centering
\caption{Test 3. Small-\(\varepsilon\) reconstruction of the trait marginal.}
\begin{tabular}{cccccccc}
\toprule
$K_W$ & $K_u$ & raw peak & basic peak & full peak & full mass & $\sigma_P$ & CPU time\\
\midrule
8 & 32 & -- & -- & -- & -- & -- & --\\
8 & 64 & -- & -- & -- & -- & -- & --\\
16 & 64 & -- & -- & -- & -- & -- & --\\
64 & 64 & -- & -- & -- & -- & -- & --\\
\bottomrule
\end{tabular}
\end{table}

The same postprocessing is also applied to direct density snapshots when they
are available, using \(\log P_k=\log\int_\Omega n(x,\theta_k)\,dx\).  This gives
a fair comparison of peak localization while keeping the WKB-specific
phase-amplitude curvature information as an additional option for the full WKB
reconstruction.

\subsection{Test 4: coarse trait grids and AP advantage of the WKB phase}\label{subsec:test_coarse_AP_new}

This test demonstrates the main computational advantage of the WKB formulation.  We keep \(\varepsilon\) small and vary the number of trait grid points:
\[
    J=200,
    \qquad
    \varepsilon=5\times10^{-3},
    \qquad
    K\in\{16,24,32,48,64,128\}.
\]
A high-resolution direct computation or, when this is too expensive, the limiting trait \(\theta_m\), is used as the reference.  We compare the selected traits
\[
    e_{\rm dir}=d_{\mathbb T}(\theta_{\rm dir},\theta_{\rm ref}),
    \qquad
    e_{\rm WKB}=d_{\mathbb T}(\theta_{\rm WKB},\theta_{\rm ref}).
\]
The expected outcome is that the direct density solver requires a much finer trait grid to localize the narrow peak, whereas the WKB phase remains a smoother object and identifies the concentration location on substantially coarser trait grids.  {\color{red} In the dual-grid version we distinguish the cost of the one-dimensional phase grid from the cost of the space-dependent amplitude grid.  Increasing \(K_u\) improves phase peak localization at a relatively small cost, while increasing \(K_W\) multiplies the number of spatial amplitude unknowns.  This distinction is minor in one spatial dimension but becomes important in two or more spatial dimensions, where the amplitude solve is the dominant cost.}

\begin{table}[ht]
\centering
\caption{Test 4. Coarse trait-grid comparison and AP advantage.}
\begin{tabular}{ccccccc}
\toprule
$K$ & $\theta_{\rm ref}$ & $\theta_{\rm dir}$ & $\theta_{\rm WKB}$ & $e_{\rm dir}$ & $e_{\rm WKB}$ & CPU ratio\\
\midrule
16 & -- & -- & -- & -- & -- & --\\
24 & -- & -- & -- & -- & -- & --\\
32 & -- & -- & -- & -- & -- & --\\
48 & -- & -- & -- & -- & -- & --\\
64 & -- & -- & -- & -- & -- & --\\
128 & -- & -- & -- & -- & -- & --\\
\bottomrule
\end{tabular}
\end{table}

\subsection{Test 5: accuracy and convergence tests}\label{subsec:test_accuracy_new}

The fifth test checks standard numerical accuracy.  Since closed-form solutions are not available, we compute a reference solution on a refined grid and compare coarser computations with it.  For example, use
\[
    \varepsilon=2\times10^{-2},
    \qquad
    (J_{\rm ref},K_{\rm ref},\tau_{\rm ref})=(400,256,2.5\times10^{-4}),
\]
and compare against coarser triples such as
\[
    (J,K,\tau)\in
    \{(100,64,10^{-3}), (200,128,5\times10^{-4}), (300,192,3.33\times10^{-4})\}.
\]
The reported quantities are the errors in \(\rho\), the trait marginal \(P_k\), the phase \(u_k\) up to the gauge \(\max_k u_k=0\), and the selected trait.  If different trait grids are used, the reference phase and marginal are interpolated periodically in \(\theta\). {\color{red} We also perform a time-step sensitivity test for the WKB evolution.  This test records \(E^n_{\rm gi}\), its phase and amplitude components, the selected trait \(\theta_{\rm WKB}\), and the diagnostics associated with the projective time integrator.  The purpose is to separate two questions.  The first one is whether the WKB phase correctly captures the fittest trait on a coarse trait grid.  The second one is whether the full amplitude and density variables have reached a fully converged steady state.  In the very small mutation regime, the first may hold before the second.}

\begin{table}[ht]
\centering
\caption{Test 5. Accuracy against a refined WKB reference solution.}
\begin{tabular}{cccccc}
\toprule
$J$ & $K$ & $\tau$ & error in $\rho$ & error in $P$ & trait error\\
\midrule
100 & 64 & $10^{-3}$ & -- & -- & --\\
200 & 128 & $5\times10^{-4}$ & -- & -- & --\\
300 & 192 & $3.33\times10^{-4}$ & -- & -- & --\\
\bottomrule
\end{tabular}
\end{table}

\subsection{Test 6: two-dimensional spatial experiment}\label{subsec:test_2d_new}

The last test checks that the WKB strategy is not restricted to one spatial dimension.  We solve the model on
\[
    \Omega=(0,1)^2
\]
using a triangular finite-element discretization in \(x=(x_1,x_2)\) and the same periodic finite-difference discretization in \(\theta\).  A representative two-dimensional landscape is
\begin{equation}\label{eq:K_2d_new}
    K(x_1,x_2)=1+0.35\cos(2\pi x_1)\cos(2\pi x_2)+0.15\cos(4\pi x_1),
\end{equation}
with the same baseline dispersal function \eqref{eq:num_D_baseline_new}.  Use
\[
    K_\theta\in\{32,64\},
    \qquad
    \varepsilon\in\{2\times10^{-2},10^{-2}\}.
\]
The outputs are the spatial total density \(\rho(x_1,x_2)\), the trait marginal \(P_k\), the phase \(u_k\), and the selected trait \(\theta_{\rm WKB}\).  This test is mainly qualitative: it should show that the spatial heterogeneity is resolved by the finite-element amplitude/eigenvalue solves, while the concentration location is still captured through the one-dimensional phase variable.

\begin{table}[ht]
\centering
\caption{Test 6. Two-dimensional spatial WKB experiment.}
\begin{tabular}{ccccc}
\toprule
$\varepsilon$ & $K_\theta$ & mesh size & $\theta_{\rm WKB}$ & $d_{\mathbb T}(\theta_{\rm WKB},\theta_m)$\\
\midrule
$2\times10^{-2}$ & 32 & -- & -- & --\\
$2\times10^{-2}$ & 64 & -- & -- & --\\
$10^{-2}$ & 32 & -- & -- & --\\
$10^{-2}$ & 64 & -- & -- & --\\
\bottomrule
\end{tabular}
\end{table}



\appendix

\section{Numerical Hamiltonians and conservative fluxes}
\label{app:hamiltonian_flux}

For the monotone theory in Section~3, the numerical Hamiltonian
\(\widehat H(a,b)\) is assumed to be consistent with \(-p^2\), nondecreasing in
its first argument, and nonincreasing in its second argument.  Two standard
choices are the Lax--Friedrichs Hamiltonian
\begin{equation}
\label{eq:app_LF_H}
    \widehat H_{\rm LF}(a,b)
    =
    -\left(\frac{a+b}{2}\right)^2
    +
    \frac{\alpha}{2}(a-b),
\end{equation}
where \(\alpha\) is chosen no smaller than the relevant characteristic speed,
and the Godunov Hamiltonian
\begin{equation}
\label{eq:app_Godunov_H}
    \widehat H_G(a,b)
    =
    -(a^-)^2-(b^+)^2,
    \qquad
    a^-:=\min\{a,0\},
    \qquad
    b^+:=\max\{b,0\}.
\end{equation}
The stationary phase estimate in
Proposition~\ref{prop:stationary_phase_estimate} uses
\(\widehat H_G\).

The conservative flux in the split amplitude equation is associated with
\[
    F=-W\partial_\theta u.
\]
On a uniform periodic trait grid, write
\[
    D_\theta\widehat F(W,u)_{j,k}
    =
    \frac{\widehat F_{j,k+1/2}-\widehat F_{j,k-1/2}}{\Delta\theta},
    \qquad
    a_{k+1/2}
    =
    -\frac{u_{k+1}-u_k}{\Delta\theta}.
\]
The upwind flux is
\begin{equation}
\label{eq:app_upwind_flux}
    \widehat F^{\rm up}_{j,k+1/2}
    =
    a_{k+1/2}^+ W_{j,k}
    +
    a_{k+1/2}^- W_{j,k+1},
    \qquad
    a^+=\max\{a,0\},
    \qquad
    a^-=\min\{a,0\}.
\end{equation}
A Lax--Friedrichs flux is
\begin{equation}
\label{eq:app_LF_flux}
    \widehat F^{\rm LF}_{j,k+1/2}
    =
    \frac12 a_{k+1/2}(W_{j,k}+W_{j,k+1})
    -
    \frac12\alpha_{k+1/2}(W_{j,k+1}-W_{j,k}),
    \qquad
    \alpha_{k+1/2}\ge |a_{k+1/2}|.
\end{equation}
Both fluxes are conservative in the periodic trait variable.  The upwind flux
also satisfies the quasi-positivity condition used in
Lemma~\ref{lem:positivity_W}.

Some numerical experiments use high-order one-sided WENO reconstructions for
the derivatives entering the Hamiltonian.  This is a numerical option for
improving resolution of the phase profile.  It is not used in the monotonicity
argument of Section~3, and the AP estimate is stated for the monotone
Hamiltonians above.


\section{Scalar mass-balance correction}
\label{app:mass_correction}

We give the algebraic form of the scalar mass-balance correction used in the
fully discrete method.  All spatial integrals below are understood in the
discrete sense and may be evaluated by any consistent spatial quadrature
associated with the fully discrete scheme.

The correction is based on the mass identity
\[
    \varepsilon \frac{d}{dt}\int_\Omega \rho(x,t)\,dx
    =
    \int_\Omega \rho(x,t)\bigl(K(x)-\rho(x,t)\bigr)\,dx .
\]
Let \(\widetilde W^{n+1}\), \(\widetilde u^{n+1}\), and
\(\widetilde\rho^{n+1}\) denote the provisional amplitude, phase, and
reconstructed spatial marginal after one time step.  The correction rescales
only the amplitude,
\[
    W^{n+1}=\lambda\widetilde W^{n+1},
    \qquad
    u^{n+1}=\widetilde u^{n+1},
    \qquad
    \lambda>0 .
\]
Consequently,
\[
    \rho^{n+1}=\lambda\widetilde\rho^{n+1},
    \qquad
    n^{n+1}=\lambda\widetilde n^{n+1},
    \qquad
    P^{n+1}=\lambda\widetilde P^{n+1}.
\]
Thus the correction changes only the density scale and leaves the phase profile
unchanged.

For compactness, define
\[
    \widetilde M
    =
    \int_\Omega \widetilde\rho^{n+1}\,dx,
    \qquad
    \widetilde S
    =
    \int_\Omega K\widetilde\rho^{n+1}\,dx,
    \qquad
    \widetilde Q
    =
    \int_\Omega (\widetilde\rho^{n+1})^2\,dx,
\]
and
\[
    M^n
    =
    \int_\Omega \rho^n\,dx,
    \qquad
    R^n
    =
    \int_\Omega \rho^n(K-\rho^n)\,dx .
\]
If the trapezoidal discretization is applied in time, \(\lambda\) is determined
by
\[
    \varepsilon
    \frac{\lambda\widetilde M-M^n}{\tau}
    =
    \frac12
    \left[
        R^n+\lambda\widetilde S-\lambda^2\widetilde Q
    \right].
\]
Equivalently,
\[
    \mathcal A_{\rm TR}\lambda^2
    +
    \mathcal B_{\rm TR}\lambda
    +
    \mathcal C_{\rm TR}
    =
    0,
\]
where
\[
    \mathcal A_{\rm TR}
    =
    \frac{\tau}{2}\widetilde Q,
    \qquad
    \mathcal B_{\rm TR}
    =
    \varepsilon\widetilde M
    -
    \frac{\tau}{2}\widetilde S,
    \qquad
    \mathcal C_{\rm TR}
    =
    -\varepsilon M^n
    -
    \frac{\tau}{2}R^n .
\]
A backward Euler version is obtained from
\[
    \varepsilon
    \frac{\lambda\widetilde M-M^n}{\tau}
    =
    \lambda\widetilde S-\lambda^2\widetilde Q,
\]
or equivalently
\[
    \mathcal A_{\rm BE}\lambda^2
    +
    \mathcal B_{\rm BE}\lambda
    +
    \mathcal C_{\rm BE}
    =
    0,
\]
with
\[
    \mathcal A_{\rm BE}=\tau\widetilde Q,
    \qquad
    \mathcal B_{\rm BE}=\varepsilon\widetilde M-\tau\widetilde S,
    \qquad
    \mathcal C_{\rm BE}=-\varepsilon M^n .
\]
The positive root is selected in both cases.  When the provisional
reconstruction already satisfies the discrete mass balance, the correction
factor is close to one.  In the small-mutation regime, this scalar projection
helps stabilize the total mass and the peak height of the trait marginal
without changing the phase-based peak location.


\section{Density and trait-marginal reconstruction}
\label{app:reconstruction}

We summarize the reconstruction procedures used to recover density-level
quantities from the WKB variables.  The basic stable reconstruction of
\(\rho\) is the log-stabilized quadrature.  Define
\[
    \ell_{j,k}
    =
    \log\bigl(\max\{W_{j,k},W_{\rm floor}\}\bigr)
    +
    \frac{u_k}{\varepsilon}.
\]
Then
\[
    \rho_j
    =
    \sum_k\omega_k\exp(\ell_{j,k})
\]
is evaluated by a row-wise log-sum-exp formula.  This is algebraically
equivalent to the composite quadrature
\[
    \rho_{j,Q}
    =
    \Delta\theta\sum_k W_{j,k}\exp(u_k/\varepsilon),
\]
but avoids underflow and overflow when \(\varepsilon\) is small.

A local Laplace reconstruction may be used near a nondegenerate maximizer of
the phase.  Let \(k_*\) be a nodal maximizer of \(u_k\), and set
\[
    d_{\theta\theta}u_{k_*}
    =
    \frac{u_{k_*+1}-2u_{k_*}+u_{k_*-1}}{(\Delta\theta)^2}.
\]
When \(d_{\theta\theta}u_{k_*}<0\), a three-point parabolic refinement gives
\[
    z_*
    =
    -\frac{\Delta\theta}{2}
    \frac{u_{k_*+1}-u_{k_*-1}}
         {u_{k_*+1}-2u_{k_*}+u_{k_*-1}},
    \qquad
    \theta_*=(\theta_{k_*}+z_*)\bmod 1 .
\]
Interpolating \(W\) and \(u\) locally at \(\theta_*\), one obtains the Laplace
approximation
\begin{equation}
\label{eq:laplace_rho}
    \rho_j^{\rm Lap}
    =
    W(x_j,\theta_*)
    \exp\left(\frac{u(\theta_*)}{\varepsilon}\right)
    \left(\frac{2\pi\varepsilon}{|d_{\theta\theta}u_{k_*}|}\right)^{1/2}.
\end{equation}
If the curvature is not negative enough or the local amplitude is not positive,
the method falls back to the log-stabilized quadrature.

The trait marginal is reconstructed in logarithmic form.  First define
\[
    A_k=\int_\Omega W(x,\theta_k)\,dx,
\]
and then
\begin{equation}
\label{eq:appendix_logP}
    L_k=\log P_k
    =
    \frac{u_k}{\varepsilon}
    +
    \log\bigl(\max\{A_k,A_{\rm floor}\}\bigr).
\end{equation}
Interpolation and peak localization are applied to \(L=\log P\), not to
\(P\).  A WENO interpolation of \(L_k\) is used on a fine periodic search grid,
followed by a local parabolic refinement to obtain the peak location.

After the peak location has been found, a local quadratic reconstruction gives
the peak width.  If
\[
    u(\theta_*+z)\simeq \alpha_0+\alpha_1z+\alpha_2z^2,
    \qquad
    \log A(\theta_*+z)\simeq \beta_0+\beta_1z+\beta_2z^2,
\]
then
\[
    \log P(\theta_*+z)
    \simeq
    \left(\frac{\alpha_0}{\varepsilon}+\beta_0\right)
    +
    \left(\frac{\alpha_1}{\varepsilon}+\beta_1\right)z
    +
    \left(\frac{\alpha_2}{\varepsilon}+\beta_2\right)z^2 .
\]
When \(\alpha_2/\varepsilon+\beta_2<0\), the fitted peak width is
\begin{equation}
\label{eq:appendix_peak_width}
    \sigma_P
    \simeq
    \frac{1}{\sqrt{-2(\alpha_2/\varepsilon+\beta_2)}} .
\end{equation}
The reconstructed shape may then be normalized by a prescribed mass \(M_{\rm
target}\),
\begin{equation}
\label{eq:appendix_mass_shape}
    P_{\rm rec}(\theta)
    =
    M_{\rm target}
    \frac{
    \exp(\widehat L(\theta)-\max_\eta\widehat L(\eta))
    }
    {
    \int_{\mathbb T}
    \exp(\widehat L(\eta)-\max_\xi\widehat L(\xi))\,d\eta
    } .
\end{equation}
This normalization stabilizes the reconstructed peak height after the location
and width have been determined.

For the dual trait-grid implementation, let \(\Theta_u\) be the fine phase grid
and \(\Theta_W\) the coarse amplitude grid.  The amplitude is first interpolated
to the phase grid,
\[
    W_U=I_{W\to u}W_c,
\]
and the log-stabilized quadrature is then performed on \(\Theta_u\):
\[
    \rho_j
    =
    \sum_{\theta_k\in\Theta_u}
    \omega_k
    \exp\left(
    \log\bigl(\max\{(W_U)_{j,k},W_{\rm floor}\}\bigr)
    +
    \frac{u_k}{\varepsilon}
    \right).
\]
Thus the phase and peak localization are resolved on the fine one-dimensional
trait grid, while the space-dependent amplitude is stored and evolved on the
coarser trait grid.


\section{A sufficient condition for the one-sided remainder bound}
\label{app:remainder}

We give a simple situation in which Assumption~\ref{ass:remainder} can be
verified.  Suppose that the reconstructed density satisfies the conservative
semi-discrete density equation
\begin{equation}
\label{eq:app_density_compatible_n}
    \varepsilon\frac{\mathrm d}{\mathrm dt}n_{j,k}^\varepsilon
    -
    D_k\delta_x^2 n_{j,k}^\varepsilon
    -
    \varepsilon^2\delta_\theta^2 n_{j,k}^\varepsilon
    =
    \bigl(K_j-\rho_{j,Q}^\varepsilon\bigr)n_{j,k}^\varepsilon .
\end{equation}
Then, in the notation of \eqref{eq:semi_discrete_n},
\[
    R_{j,k}^\varepsilon
    =
    \varepsilon^2\delta_\theta^2 n_{j,k}^\varepsilon .
\]
Therefore
\[
    R_j^{\varepsilon,D}
    =
    \varepsilon^2\Delta\theta
    \sum_k
    \frac{\delta_\theta^2 n_{j,k}^\varepsilon}{D_k}.
\]
Periodic summation by parts gives
\[
    R_j^{\varepsilon,D}
    =
    \varepsilon^2\Delta\theta
    \sum_k
    n_{j,k}^\varepsilon
    \delta_\theta^2\left(\frac{1}{D_k}\right).
\]
Hence, for \(0<\varepsilon\le 1\),
\[
    \bigl(R_j^{\varepsilon,D}\bigr)_+
    \le
    C_{D,h}m_j^\varepsilon,
    \qquad
    C_{D,h}
    =
    \max_k
    \left|
    \delta_\theta^2\left(\frac{1}{D_k}\right)
    \right|,
    \qquad
    m_j^\varepsilon
    =
    \Delta\theta\sum_k n_{j,k}^\varepsilon .
\]
This verifies Assumption~\ref{ass:remainder} with \(C_R=C_{D,h}\) and
\(r_h=0\).  If \(D\) is smooth and the standard centered periodic difference is
used, \(C_{D,h}\) remains bounded under regular trait-grid refinement.


\section{Direct density solver and two-dimensional finite-element realization}
\label{app:direct_2d}

For comparison, the original density equation can be discretized by the linear
implicit update
\begin{equation}
\label{eq:direct_density_fully}
    \varepsilon
    \frac{n_{j,k}^{n+1}-n_{j,k}^n}{\tau}
    -
    D_k\delta_x^2n_{j,k}^{n+1}
    -
    \varepsilon^2\delta_\theta^2n_{j,k}^{n+1}
    =
    n_{j,k}^{n+1}(K_j-\rho_j^n),
    \qquad
    \rho_j^n=\Delta\theta\sum_k n_{j,k}^n.
\end{equation}
A positivity-oriented Patankar variant is
\begin{equation}
\label{eq:direct_density_patankar}
    \varepsilon
    \frac{n_{j,k}^{n+1}-n_{j,k}^n}{\tau}
    -
    D_k\delta_x^2n_{j,k}^{n+1}
    -
    \varepsilon^2\delta_\theta^2n_{j,k}^{n+1}
    =
    (K_j-\rho_j^n)_+n_{j,k}^n
    -
    (K_j-\rho_j^n)_-n_{j,k}^{n+1}.
\end{equation}
The same spatial and trait grids are used for the direct and WKB solvers in the
one-dimensional comparison tests.

For the two-dimensional spatial experiment, the spatial Laplacian and the
eigenvalue problem are discretized by linear finite elements on a triangular
mesh.  For each trait node \(\theta_k\), the discrete effective Hamiltonian is
computed from
\begin{equation}
\label{eq:fem_H}
    A_k(\rho_Q)N^{(k)}
    =
    H_k(\rho_Q)M N^{(k)},
\end{equation}
where \(M\) is the finite-element mass matrix and
\begin{equation}
\label{eq:fem_A}
    (A_k)_{ab}
    =
    D_k\int_\Omega \nabla\varphi_a\cdot\nabla\varphi_b\,dx
    -
    \int_\Omega (K(x)-\rho_Q(x))\varphi_a\varphi_b\,dx.
\end{equation}
The Neumann boundary condition is natural in this formulation.  The phase
equation remains one-dimensional in \(\theta\), while the amplitude equation is
solved for all spatial finite-element degrees of freedom and trait nodes.





\bibliographystyle{plain}
\bibliography{rml_wkb_references}
%\begin{thebibliography}{99}
%
%\bibitem{AlmeidaPerthameRuan}
%L. Almeida, B. Perthame, and X. Ruan,
%An asymptotic preserving scheme for capturing concentrations in
%age-structured models arising in adaptive dynamics,
%\emph{Journal of Computational Physics}, 464 (2022), 111335.
%
%\bibitem{BarlesPerthame}
%G. Barles and B. Perthame,
%Concentrations and constrained Hamilton--Jacobi equations arising in adaptive
%dynamics,
%\emph{Contemporary Mathematics}, 439 (2007), 57--68.
%
%\bibitem{MirrahimiPerthame}
%S. Mirrahimi and B. Perthame,
%Asymptotic analysis of a selection model with space,
%\emph{Journal de Math\'{e}matiques Pures et Appliqu\'{e}es}, 104 (2015),
%1108--1118.
%
%\bibitem{PerthameSouganidis}
%B. Perthame and P. E. Souganidis,
%Rare mutations limit of a steady state dispersal evolution model,
%\emph{Mathematical Modelling of Natural Phenomena}, 11 (2016), 154--166.
%
%\end{thebibliography}

\end{document}
